\section{Literature}\label{sec:literature}

The study of two-class priority queues with jockeying and abandonment sits at the intersection of three well-developed streams: the exact generating-function analysis of non-preemptive priority queues, the theory of queues with customer impatience (abandonment), and the theory of jockeying and dynamic priority changes. The applied motivation throughout is the management of access to care ---elderly-care waiting lists, emergency departments, and transplant registries--- in which the two mechanisms studied here carry direct clinical readings: a waiting patient whose condition deteriorates changes priority class (jockeying), while one who leaves the list by death or transfer abandons the system (abandonment). This thesis contributes an exact, single-server treatment that admits jockeying and abandonment within a unified bivariate-PGF framework, filling a gap that the multi-server and approximate literature has so far left open.

% ── §2.1 ──────────────────────────────────────────────────────────────────────
\subsection*{Non-preemptive priority queues: exact queue-length analysis}

The expected waiting times of the non-preemptive priority $M/M/1$ queue are standard textbook material~\cite{adan2002queueing}; the exact \emph{queue-length} analysis is where the generating-function tradition begins. The direct methodological ancestor of this work is Cohen~\cite{cohen1956delay}, who studied a single server loaded by two independent Poisson streams under non-preemptive priority and introduced a partial-generating-function recursion that resolves the lower-priority queue length by transforming only one coordinate of the joint distribution while keeping the other discrete; we refer to this device as the ``Cohen trick''. It is exactly the device reused, in modern notation, in the partial-PGF analysis of Model-$A$ on the total-count state space (\S\ref{ssec:A:Stilde}). The stationary queue-length law of the two-class non-preemptive $M/M/1$ priority queue ---the content of Theorem~\ref{thm:A:PGF} (Model-$A$)--- is in this sense classical, recoverable by several routes.

The most recent exact closed-form treatment of the non-preemptive priority queue is that of Zuk \& Kirszenblat~\cite{zuk2023explicit}, who obtain the joint and marginal queue-length distributions of a non-preemptive $M/M/c$ queue with two priority classes and a common service rate by generating-function methods and a quadratic recurrence; their single-server specialisation ($c=1$) coincides with Model-$A$.

% ── §2.2 ──────────────────────────────────────────────────────────────────────
\subsection*{Priority queues with abandonment}

Abandonment enters queueing theory through the Erlang-A ($M/M/c+M$) model, in which each waiting customer holds an independent exponential patience clock; the modern treatment and its heavy-traffic scaling are due to Garnett et al.~\cite{garnett2002designing}. The analytic signature of impatience is well documented: per-customer patience clocks make the transition rates state-inhomogeneous, so the generating-function solutions, where they close at all, close in terms of hypergeometric series.

Kapodistria~\cite{kapodistria2011synchronized} exhibits this pattern most clearly in a single-server queue with \emph{synchronised} abandonments, in which all waiting customers abandon simultaneously, each with probability $p$, at Poisson opportunity epochs. The stationary law is solved exactly through basic ($q$-)hypergeometric series, and in the no-synchronisation limit $p\to0^+$ it recovers the $M/M/1+M$ queue-length PGF as a ratio of Kummer functions: precisely the special-function structure that carries the boundary function of Model-$C_2$.

Exact stationary results for priority queues \emph{with} abandonment are scarce, and those available concern multi-server systems. Jouini \& Roubos~\cite{jouini2014multiple} treat an $M/M/s+M$ queue with two non-preemptive classes in which the classes are statistically identical ---a common service rate and a common patience distribution, so that the classes differ only in their priority level. They derive Laplace--Stieltjes transforms for conditional and unconditional waiting times together with the probabilities of service and of abandonment; this is the multi-server analogue of Model-$C_2$, but with the waiting time, not the queue-length PGF, as its object. Where exact transforms are unavailable the literature turns to fluid and diffusion approximations or to many-server asymptotics~\cite{zeltyn2005callcenters}; the exact Erlang-A stationary law itself is expressed through confluent hypergeometric and incomplete-gamma functions, the same special-function structure that appears, at finite scale, in the closed form for Model-$C_2$.

The exact multi-server priority results that do exist involve no abandonment at all: Wang et al.~\cite{wang2015mmc} analyse a preemptive $M/M/c$ queue with class-dependent service rates and no impatience, and Zuk \& Kirszenblat~\cite{zuk2023explicit} a non-preemptive $M/M/c$ queue, likewise without impatience. Model-$C_2$ therefore appears to be the first exact, single-server, queue-length treatment of \emph{class-asymmetric} abandonment, with impatience acting on the priority class alone. The clinical reading that licenses this asymmetry is long-standing: Zenios~\cite{zenios1999modeling} models the organ-transplant waiting list as a queue whose abandonment mechanism is patient mortality, precisely the regime in which losing high-priority demand is the outcome of interest.

% ── §2.3 ──────────────────────────────────────────────────────────────────────
\subsection*{Jockeying and priority changes}

The closest predecessor on the jockeying side is de Waal~\cite{deWaal1992approximate}, who studies an $M/G/1$ priority queue in which impatient customers \emph{may change priority class} while waiting. de Waal's analysis is approximate by necessity rather than by choice: with general service times an exact analysis is out of reach, and the paper delivers Laplace--Stieltjes-based \emph{approximations} of the waiting-time distributions rather than an exact stationary law. The healthcare incarnation of the same idea is the model of Hu, Chan \& Dong~\cite{huchandong2022proactive}, who let customers transition between a moderate and an urgent class as their condition deteriorates or improves and characterise a near-optimal $c\mu/\theta$-type scheduling index in heavy traffic; the system is multi-server and its results are asymptotic, so it too is approximate in the sense of Table~\ref{tab:literature}.

A complementary strand studies dynamic rather than static priority: Stanford et al.~\cite{stanford2014accumulating} analyse the \emph{accumulating-priority} $M/G/1$ queue, in which each class earns priority linearly in its waiting time so that older customers eventually overtake newer ones, and derive exact waiting-time distributions through a maximum-priority process and Laplace--Stieltjes transforms. Accumulating priority is a continuous-time analogue of the discrete, rate-$\gamma_i$ jockeying studied here; like the present work it is exact and single-server, but its object is the sojourn time rather than the bivariate queue-length PGF.

Model-$B_2$ is the exact counterpart of this ``priority-change'' literature. It retains the deterioration mechanism ---a waiting class-$1$ customer relocates to the class-$2$ queue at rate $\gamma_1$--- and pays for exactness with the single-server, Markovian, one-directional restriction. In return it produces a closed-form boundary function in place of an approximation.

% ── §2.4 ──────────────────────────────────────────────────────────────────────
\subsection*{Methodological lineage}

The functional-equation technique used in Sections~\ref{sec:model_b2} and~\ref{sec:model_c2} belongs to the kernel-method and boundary-value tradition for two-dimensional Markovian queues. The systematic theory ---reducing a bivariate balance equation to the curve on which its kernel vanishes, and from there to a boundary-value problem for the unknown boundary functions--- is developed by Fayolle, Iasnogorodski \& Malyshev~\cite{fayolle1999random} for random walks in the quarter-plane. That approach still delivers explicit two-class results: Devos, Walraevens, Fiems \& Bruneel~\cite{devos2020alternating} obtain the closed-form joint queue-length distribution of a discrete-time two-class queue with randomly alternating service by analysing the kernel of just such a functional equation and determining its two unknown boundary functions through analytic continuation.

The arguments of this thesis are elementary, one-dimensional instances of that method. In Model-$B_2$ the boundary function is fixed by demanding \emph{analyticity} of the PGF at the interior point where the kernel root meets the diagonal (\S\ref{sec:model_b2}), and in Model-$C_2$ by demanding mere \emph{boundedness} at the unit-disc boundary (\S\ref{sec:model_c2}); both select the analytic solution of the reduced equation, exactly the role the boundary conditions play in the general theory.

Two narrower tools recur. The confluent hypergeometric (Kummer) function that carries the Model-$C_2$ boundary function is the same special function through which the Erlang-A and reneging-queue stationary laws are expressed~\cite{zeltyn2005callcenters,kapodistria2011synchronized,dlmf}; and the selection of the decaying solution of the three-term recurrence behind the partial-PGF analysis is governed by Pincherle's theorem, for which we cite the recurrence-relation treatment of Gautschi~\cite{gautschi1967computational} alongside the DLMF~\cite{dlmf}.

% ── §2.6 ──────────────────────────────────────────────────────────────────────
\subsection*{Position of this thesis}

\begin{table}[H]
\centering
\caption{Literature overview: models and mechanisms.
\checkmark~= present, --~= absent, $\approx$~= approximate only.}
\label{tab:literature}
\renewcommand{\arraystretch}{1.3}
\footnotesize
\begin{tabular}{lcccccc}
\hline
\textbf{Reference} & \textbf{Servers} & \textbf{Non-preem.} & \textbf{Jockeying}
  & \textbf{Aband.} & \textbf{Queue-len.} & \textbf{Exact} \\
\hline
Cohen (1956)                      & 1   & \checkmark & -- & --         & \checkmark & \checkmark \\
de Waal (1992)                    & 1   & \checkmark & \checkmark & --  & --         & $\approx$ \\
Jouini \& Roubos (2014)           & $c$ & \checkmark & -- & \checkmark & --         & \checkmark \\
Wang et al.\ (2015)               & $c$ & --         & -- & --         & \checkmark & \checkmark \\
Stanford et al.\ (2014)           & 1   & \checkmark & \checkmark & --  & --         & \checkmark \\
Zuk \& Kirszenblat (2023)         & $c$ & \checkmark & -- & --         & \checkmark & \checkmark \\
Hu, Chan \& Dong (2022)           & $c$ & \checkmark & \checkmark & \checkmark & --   & $\approx$ \\
\hline
\textbf{This thesis (Model-$A$)}    & \textbf{1} & \checkmark & -- & --  & \checkmark & \checkmark \\
\textbf{This thesis (Model-$B_2$)}& \textbf{1} & \checkmark & \checkmark & -- & \checkmark & \checkmark \\
\textbf{This thesis (Model-$C_2$)}& \textbf{1} & \checkmark & -- & \checkmark & \checkmark & \checkmark \\
\hline
\end{tabular}
\end{table}

Table~\ref{tab:literature} places this work against its closest neighbours, and reading it by row makes the gap precise. The Cohen (1956) row is the single-server, exact, non-preemptive queue-length baseline with neither extra mechanism ---the ancestor of Model-$A$. de Waal (1992) adds jockeying (priority change) at a single server, but its analysis is approximate and concerns waiting times rather than queue lengths. Jouini \& Roubos (2014) add abandonment exactly, but at $c$ servers, for statistically identical classes, and again for waiting times. Wang et al.\ (2015) is exact and queue-length-focused, yet multi-server, preemptive, and mechanism-free. Stanford et al.\ (2014) is single-server and exact with a dynamic-priority (jockeying-like) mechanism ---the closest analogue of Model-$B_2$ in the table--- but its object is the sojourn time. Zuk \& Kirszenblat (2023) is the exact multi-server queue-length result whose $c=1$ case is Model-$A$, with neither jockeying nor abandonment. Hu, Chan \& Dong (2022) carries both mechanisms ---jockeying-as-deterioration and abandonment--- and is in that sense the richest model in the table, but it is multi-server and offers no exact results. The pattern is symptomatic of the multi-server literature as a whole: once $c>1$ is combined with jockeying or abandonment, the analysis moves to approximations and asymptotics, so that exact, mechanism-rich analysis is essentially a single-server phenomenon.

No row above combines the exact bivariate queue-length PGF with either jockeying or abandonment ---which is precisely the column pattern of the three ``This thesis'' rows. The gap can therefore be stated in one sentence: prior to this work there is no exact, single-server, queue-length analysis that places non-preemptive priority together with jockeying or abandonment inside a single PGF framework.

The limits of these claims must be stated with equal precision. The framework is unified but not complete: the full model (Model-$X$, Remark~\ref{rem:B:modelX}), in which jockeying and abandonment act together, and the bidirectional-jockeying model (Model-$B$, \S\ref{sec:model_b}) are left open. For each of these, this thesis characterises the analytical obstruction and reduces the problem to a Volterra-type integral equation, providing partial results rather than a closed form. The exact contributions are confined to the three ``This thesis'' rows of Table~\ref{tab:literature} ---the baseline Model-$A$, the one-way jockeying Model-$B_2$, and the class-1 abandonment Model-$C_2$, together with their head-of-line simplifications. The claim is exactness and unification at a single server with one mechanism active at a time; it is explicitly not a solution of the combined or bidirectional problems, which remain, as in the multi-server literature, beyond exact reach.
